\chapter{Reading a Classic: Shuffling Cards}
\label{ch:shuffling}

\begin{center}\itshape
Some papers are read once and filed. Others are read for forty years,\\
because each generation finds in them a question the authors left open.\\
This chapter reads one of those, and then answers one of its questions.
\end{center}

In the late 1950s a fourteen-year-old New Yorker named Persi Diaconis ran away from home to go on the road with Dai Vernon, the most
admired sleight-of-hand artist of his time. He spent the next ten years as a travelling magician, learned that a deck of cards is a
small universe of $52!$ states, and eventually wanted to know how fast a shuffle explores it. He went back to school, studied
mathematics at the City College of New York, and was admitted to the graduate statistics programme at Harvard on the strength of a letter
from Martin Gardner, who had admired two of his card tricks. By the late 1970s he was a professor at Stanford, and the question ``how many
shuffles?'' had become a branch of probability theory.

The first landmark of that branch is a paper Diaconis wrote in 1981 with Mehrdad Shahshahani, \emph{Generating a random permutation with random
transpositions}. Its result, the cutoff for random transpositions, is proved in Volume~11 (Theorem~14.8). So in this chapter we do something
different from Chapter~\ref{ch:reading}: we read a paper whose content we already know. The question is no longer ``is it right?'' but ``what is
left?''. Reading a classic is mostly a matter of locating it: finding where each of its ideas now lives, which of its questions have been
answered, and which have not. We shall make a map of the paper, find in it a question that the series never answered, and answer it: the
theorem of Bayer and Diaconis that turned ``seven shuffles'' into a precise statement.

\section{The first pass}

The paper is twenty-one pages in a probability journal. Its setting is a deck of $n$ cards, shuffled by swapping two cards chosen independently and
uniformly (the same card twice means doing nothing). After $t$ swaps, let $d(t)$ be the total variation distance between the law of the deck and
the uniform law. The claims, in our notation:
\begin{itemize}
\item if $t=\frac12n\ln n+cn$ with $c>0$, then $d(t)\le ae^{-2c}$, for a universal constant $a$;
\item if $t=\frac12n\ln n-cn$ with $c>0$, then $d(t)$ is close to $1$ when $c$ is large.
\end{itemize}
So the distance drops from near $1$ to near $0$ in a window of width of order $n$ around $\frac12n\ln n$: the \emph{cutoff phenomenon}. The method is
Fourier analysis on the symmetric group. The authors bound $d(t)$ by a sum over the irreducible
representations of $\mathfrak S_n$, and bound that sum with the characters of $\mathfrak S_n$ at a transposition, which Frobenius had computed in 1900.

\emph{Five questions.} Claim: cutoff at $\frac12n\ln n$, window $O(n)$. Novelty: the first sharp analysis of a natural random walk on a large
non-abelian group. Type: new method. Tools: representation theory, the upper bound lemma (proved in the paper), Frobenius's formula. Worth a second
pass? Certainly, for the method, which became standard.

\section{The second pass: a reading map}

For a classic, the second pass is best done as a map: for each idea of the paper, where it is proved now. Here is ours, idea by idea.

\begin{itemize}
\item Fourier transform and Plancherel on a finite group: Volume~11, Theorem~14.2.
\item The upper bound lemma: Volume~11, Lemma~14.4.
\item For a class function, $\hat Q(\rho)$ is a scalar: Volume~11, Proposition~14.3.
\item The scalar is a character ratio $s(\lambda)$: Volume~11, Proposition~14.5.
\item Frobenius: $\chi^\lambda(\tau)/f^\lambda$ is a sum of contents: Volume~10, Theorem~10.15.
\item Bounds on $s(\lambda)$, counting $\sum(f^\lambda)^2$: Volume~11, Lemmas~14.6 and~14.7.
\item Upper bound at $\frac12n\ln n+cn$: Volume~11, Theorem~14.8, with $3e^{-4c/3}$ for $n\ge10^4$.
\item Lower bound at $\frac12n\ln n-cn$: Volume~4, Proposition~8.21 (untouched cards).
\item Cutoff: Volume~11, Corollary~14.9.
\end{itemize}

The map shows what we have and what we lack. We have every idea of the paper. We lack two things: the sharp constant (the paper's $ae^{-2c}$
against our $3e^{-4c/3}$), and any statement about the \emph{shape} of the drop. Both are worth understanding, and the second is where research went
next.

\emph{The constant.} The upper bound lemma cannot give better than $e^{-2c}$: the single term of $\lambda=(n-1,1)$ in the sum is
$(n-1)^2(1-\frac2n)^{2t}$, which at $t=\frac12n\ln n+cn$ is about $e^{-4c}$ (Exercise~\ref{ex:lemmalimit}). So the paper's bound is the best its method can give,
and our $e^{-4c/3}$ comes only from a crude treatment of the shapes with a long first row. The third pass of the paper would be to redo
Lemma~10.6 with the exact values of $s(\lambda)$ for $\lambda_1=n-k$ and $k$ small; it is careful bookkeeping, not a new idea, and we leave it.

\emph{The shape.} Rescale time as $t=\frac12n\ln n+cn$. Does $d(t)$ converge, for each fixed $c$, to a function $\Psi(c)$, the \emph{limit profile}? The
paper does not say. A hint comes from fixed points. Volume~11, Problem~14.4, computes the mean number of fixed points after $t$ swaps exactly:
\[
\mathbb E[F]=1+(n-1)\Bigl(1-\frac2n\Bigr)^t\longrightarrow1+e^{-2c},
\]
while a uniform permutation has on average $1$ fixed point, and in the limit a Poisson$(1)$ number of them.

\begin{remark}[the limit profile, admitted]\label{rem:teyssier}
In 2020 Lucas Teyssier proved that for random transpositions
\[
d\bigl(\lfloor\tfrac12n\ln n+cn\rfloor\bigr)\longrightarrow\bigl\|\mathrm{Poisson}(1+e^{-2c})-\mathrm{Poisson}(1)\bigr\|\qquad(n\to\infty),
\]
for every real $c$: in the limit, the only visible trace of the shuffle is the number of fixed points, which is Poisson with an excess mean $e^{-2c}$. His proof
refines the upper bound lemma itself, and uses the characters of $\mathfrak S_n$ in a much more precise way than we can here.
\end{remark}

Forty years separate the paper from the answer to its most natural question. That is typical, and it is the most important thing a reading map
can show you: a question that a classic \emph{could} have asked, did not, and that was still open decades later.

\begin{think}
Why should fixed points be the only trace of the shuffle? After $\frac12n\ln n+cn$ swaps, about how many cards have never been touched? Where are
the touched cards likely to be?
\end{think}

\section{The question we answer: riffle shuffles}

Nobody shuffles by random transpositions. People riffle: cut the deck in two and interleave the halves. Volume~4 (Definition~8.22 and
Theorem~8.23) describes the Gilbert--Shannon--Reeds model of the riffle and its generalisation, the $a$-shuffle, and proves a
remarkable exact formula: an $a$-shuffle of the ordered deck produces an arrangement $\sigma$ with probability
\[
Q_a(\sigma)=\frac1{a^n}\binom{a+n-r}n,\qquad r=r(\sigma)\text{ the number of rising sequences of }\sigma,
\]
and $k$ riffle shuffles are a $2^k$-shuffle (Volume~4, Proposition~8.24). Volume~4 then computes the distance to uniform for $52$ cards ($0.614$ after six
shuffles, $0.334$ after seven), and \emph{states} that for large $n$ the cutoff is at $\frac32\log_2n$ shuffles. It does not prove it. Here is the
theorem, as Dave Bayer and Persi Diaconis proved it in 1992 in a paper with the best title in probability, \emph{Trailing the dovetail shuffle
to its lair}.

\begin{theorem}[Bayer and Diaconis, 1992]\label{thm:bd}
Let $a_n$ be positive integers with $a_n/n^{3/2}\to\gamma\in(0,\infty)$. Then the total variation distance between an $a_n$-shuffle of $n$ cards and the
uniform law satisfies
\[
d_n\longrightarrow1-2\,\Phi\Bigl(-\frac1{4\sqrt3\,\gamma}\Bigr),
\]
where $\Phi$ is the standard normal distribution function.
\end{theorem}

For $k$ riffle shuffles, $a=2^k$, and $\gamma=2^\theta$ when $k=\frac32\log_2n+\theta$. Figure~\ref{fig:riffle15} compares the limit with the exact distances; already for $52$
cards the agreement is within $0.02$. The proof needs three facts: the distance depends only on the number of rising sequences; under the
uniform law that number is a sum of independent uniforms, rounded down; and the formula for $Q_a$ is, after taking logarithms, nearly linear in $r$.

\begin{figure}[ht]
\centering
\fig{ch03-riffle}
\caption{Distance to uniform after an $a$-shuffle with $a=2^\theta n^{3/2}$ (rounded to an integer), computed exactly from the Bayer--Diaconis formula for $n=52$,
$1000$ and $10000$ cards, against the limit $1-2\Phi(-2^{-\theta}/(4\sqrt3))$ of Theorem~\ref{thm:bd}. For $k$ riffle shuffles, $\theta=k-\frac32\log_2n$; for $52$ cards,
seven shuffles is $\theta\approx-1.55$.}
\label{fig:riffle15}
\end{figure}

\subsection*{Step 1: a single statistic}

Write $\mathbb U$ for the uniform law on arrangements and $R_a(r)=n!\,Q_a(\sigma)$ when $\sigma$ has $r$ rising sequences. Then
\[
d=\sum_\sigma\Bigl(\frac1{n!}-Q_a(\sigma)\Bigr)^+=\mathbb E_{\mathbb U}\bigl[(1-R_a(r))^+\bigr],
\]
since the total variation distance is the sum of the positive parts of the differences (both laws have mass $1$). So everything depends on the law of $r$
under the uniform law, and on the function $R_a$. Reading lesson: \emph{when the law you study is a function of one statistic, the whole problem
is about that statistic}.

\subsection*{Step 2: rising sequences and sums of uniforms}

The number of rising sequences is $r(\sigma)=1+D$, where $D$ is the number of $i<n$ such that card $i+1$ lies above card $i$. If $\pi(i)$ is the position of card $i$,
then $D$ is the number of \emph{descents} $\pi(i+1)<\pi(i)$ of the sequence $\pi(1),\dots,\pi(n)$, and under the uniform law $\pi$ is a uniform permutation.

\begin{lemma}[Tanny, 1973]\label{lem:tanny}
Let $U_1,\dots,U_n$ be independent and uniform on $[0,1]$, and $S_n=U_1+\dots+U_n$. Then the number of descents of a uniform permutation of $n$ elements
has the law of $\lfloor S_n\rfloor$.
\end{lemma}

\begin{proof}
Let $V_k$ be the fractional part of $S_k$. Then $V_1=U_1$, and $V_{k+1}$ is the fractional part of $V_k+U_{k+1}$, which is uniform on $[0,1)$ and independent of $V_1,\dots,V_k$, whatever their values
(adding an independent uniform modulo $1$ gives a uniform). So $V_1,\dots,V_n$ are independent and uniform; almost surely they are distinct, and the permutation
recording their relative order is uniform, by symmetry. Its descents are the $k$ with $V_{k+1}<V_k$. Now if $V_k+U_{k+1}<1$, then $V_{k+1}=V_k+U_{k+1}>V_k$ (almost surely) and
$\lfloor S_{k+1}\rfloor=\lfloor S_k\rfloor$; if $V_k+U_{k+1}\ge1$, then $V_{k+1}=V_k+U_{k+1}-1<V_k$ and $\lfloor S_{k+1}\rfloor=\lfloor S_k\rfloor+1$. So each descent is a carry, and the number of descents is
$\lfloor S_n\rfloor-\lfloor S_1\rfloor=\lfloor S_n\rfloor$.
\end{proof}

So $r-1$ has the law of $\lfloor S_n\rfloor$, which has mean about $\frac n2$ and standard deviation about $\sqrt{n/12}$, and is approximately normal by the central limit
theorem. Let
\[
h=r-\frac{n+1}2,\qquad Z_n=\frac h{\sqrt{n/12}}.
\]

\begin{lemma}\label{lem:risingclt}
Under the uniform law, $Z_n$ converges in law to $N(0,1)$, $\mathbb E[Z_n^2]\le3$ for $n\ge6$, and $P(|h|>n^{3/5})\to0$.
\end{lemma}

\begin{proof}
By Lemma~\ref{lem:tanny}, $h$ has the law of $(S_n-\frac n2)+\delta_n$ with $\delta_n=\lfloor S_n\rfloor-S_n+\frac12\in(-\frac12,\frac12]$. The $U_i-\frac12$ are independent, centred, bounded, with variance
$\frac1{12}$, so $W_n=(S_n-\frac n2)/\sqrt{n/12}$ converges in law to $N(0,1)$ by Lindeberg's theorem (Volume~4, Theorem~6.10; the Lindeberg condition holds trivially for
bounded summands). Since $|Z_n-W_n|\le\varepsilon_n=\frac12\sqrt{12/n}\to0$, for every $t$ and $\varepsilon>0$ and $n$ large, $P(W_n\le t-\varepsilon)\le P(Z_n\le t)\le P(W_n\le t+\varepsilon)$, and letting
$n\to\infty$ then $\varepsilon\to0$ gives $P(Z_n\le t)\to\Phi(t)$. Next, $\mathbb E[Z_n^2]\le2\mathbb E[W_n^2]+2\varepsilon_n^2=2+\frac6n\le3$ for $n\ge6$. (The sharper $\mathbb E[h^2]=\frac{n+1}{12}$ is Volume~4, Exercise~8.8.) Finally, by Chebyshev's inequality,
\[
P(|h|>n^{3/5})\le P\bigl(|S_n-\tfrac n2|>n^{3/5}-\tfrac12\bigr)\le\frac{n/12}{(n^{3/5}-\frac12)^2}\longrightarrow0.
\]
\end{proof}

\subsection*{Step 3: the formula, linearised}

Expanding the binomial coefficient,
\[
R_a(r)=\frac{n!}{a^n}\binom{a+n-r}n=\prod_{j=1}^n\frac{a-r+j}a=\prod_{j=1}^n\Bigl(1+\frac{\ell_j-h}a\Bigr),\qquad\ell_j=j-\frac{n+1}2,
\]
with the convention that $R_a(r)=0$ if some factor vanishes. The numbers $\ell_j$ are symmetric about $0$, so $\sum_j\ell_j=0$, and $\sum_j\ell_j^2=\frac{n^3-n}{12}$.

\begin{lemma}\label{lem:linearise}
Let $s=\frac{n}{a}\sqrt{\frac n{12}}$, and suppose that $n^{3/5}\le\frac n2$ and $\frac na\le\frac12$. Then for $|h|\le n^{3/5}$,
\[
\Bigl|\ln R_a(r)+sZ_n+\frac{s^2}2\Bigr|\le\frac{n^4}{a^3}+\frac{n^{11/5}}{2a^2}+\frac{s^2}{2n^2}.
\]
\end{lemma}

\begin{proof}
Let $x_j=(\ell_j-h)/a$. Then $|x_j|\le(\frac n2+n^{3/5})/a\le\frac na\le\frac12$, so every factor is positive and $\ln(1+x_j)=x_j-\frac{x_j^2}2+\rho_j$ with
$|\rho_j|\le\sum_{m\ge3}\frac{|x_j|^m}m\le\frac{|x_j|^3}{3(1-|x_j|)}\le|x_j|^3$. Now
\[
\sum_jx_j=-\frac{nh}a=-sZ_n,\qquad\sum_jx_j^2=\frac{n^3-n}{12a^2}+\frac{nh^2}{a^2}=s^2\Bigl(1-\frac1{n^2}\Bigr)+\frac{nh^2}{a^2},
\]
(the cross term vanishes because $\sum\ell_j=0$), and $\sum_j|x_j|^3\le n(\frac na)^3$. With $h^2\le n^{6/5}$ this gives the bound.
\end{proof}

\subsection*{Step 4: the limit}

\begin{proof}[Proof of Theorem~\ref{thm:bd}]
Let $s_n=\frac n{a_n}\sqrt{\frac n{12}}$; then $s_n\to s=\frac1{\gamma\sqrt{12}}=\frac1{2\sqrt3\gamma}$. Let $g(u)=(1-e^u)^+$; it takes values in $[0,1]$ and is $1$-Lipschitz (its derivative is $-e^u\in[-1,0]$ for $u<0$,
and $0$ for $u>0$). By Step~1, $d_n=\mathbb E[g(\ln R_{a_n}(r))]$, with $g(\ln0)=1$. Compare, in three steps, with $D_n=\mathbb E[g(-s_nZ_n-\frac{s_n^2}2)]$, $D'_n=\mathbb E[g(-sZ_n-\frac{s^2}2)]$ and the
limit.

\emph{From $d_n$ to $D_n$.} The right side of Lemma~\ref{lem:linearise} is $O(n^{-1/2})$, since $a_n\asymp n^{3/2}$. On $\{|h|\le n^{3/5}\}$ the two integrands differ by at most this, and
elsewhere by at most $1$. So $|d_n-D_n|\le O(n^{-1/2})+P(|h|>n^{3/5})\to0$ by Lemma~\ref{lem:risingclt}.

\emph{From $D_n$ to $D'_n$.} By the Lipschitz property, $|D_n-D'_n|\le|s_n-s|\,\mathbb E|Z_n|+\frac12|s_n^2-s^2|\to0$, as $\mathbb E|Z_n|\le\sqrt{\mathbb E[Z_n^2]}\le\sqrt3$.

\emph{From $D'_n$ to the limit.} The function $G(z)=g(-sz-\frac{s^2}2)$ is continuous and nondecreasing in $z$, with values in $[0,1)$, and strictly increasing where it is positive.
For $u\in(0,1)$ let $z_u$ be the point where $G(z_u)=u$. Then $\{G(Z_n)>u\}=\{Z_n>z_u\}$, so
\[
D'_n=\int_0^1P(G(Z_n)>u)\dd u=\int_0^1P(Z_n>z_u)\dd u\longrightarrow\int_0^1\bigl(1-\Phi(z_u)\bigr)\dd u=\mathbb E[G(Z)],
\]
with $Z\sim N(0,1)$, by Lemma~\ref{lem:risingclt} and dominated convergence (the integrands are bounded by $1$).

\emph{Computing the limit.} With $\varphi$ the normal density, $e^{-sz-s^2/2}\varphi(z)=\varphi(z+s)$, so
\[
\mathbb E[G(Z)]=\int_\R\bigl(\varphi(z)-\varphi(z+s)\bigr)^+\dd z.
\]
The integrand is positive exactly when $|z|<|z+s|$, that is $z>-\frac s2$. So the integral is $(1-\Phi(-\frac s2))-(1-\Phi(\frac s2))=1-2\Phi(-\frac s2)$, and $\frac s2=\frac1{4\sqrt3\gamma}$.
\end{proof}

The last line has a meaning worth stating: $1-2\Phi(-\frac s2)$ is the total variation distance between $N(0,1)$ and $N(-s,1)$. In the limit, \emph{the shuffled
deck is indistinguishable from a uniform one except through its number of rising sequences, which is normal and shifted by $s$ standard
deviations}. A shuffled deck has slightly too few rising sequences, and that deficit is all there is to see.

\begin{corollary}[cutoff at $\frac32\log_2n$]\label{cor:riffle}
Let $d_n(k)$ be the distance to uniform after $k$ riffle shuffles of $n$ cards. If $k_n-\frac32\log_2n\to+\infty$ then $d_n(k_n)\to0$, and if $k_n-\frac32\log_2n\to-\infty$ then $d_n(k_n)\to1$.
\end{corollary}

\begin{proof}
Write $L(\gamma)=1-2\Phi(-\frac1{4\sqrt3\gamma})$, which decreases from $1$ to $0$ as $\gamma$ goes from $0$ to $\infty$. Riffle shuffling is a Markov chain on arrangements for
which the uniform law is stationary (it is a random walk on the group $\mathfrak S_n$, and $\sum_\sigma Q_2(\sigma^{-1}\tau)=1$ for every $\tau$), so $d_n(k)$ is
nonincreasing in $k$ (Volume~4, Lemma~8.10).

Let $\varepsilon>0$, and choose $\theta$ with $L(2^\theta)<\varepsilon$. Let $k'_n=\lceil\frac32\log_2n+\theta\rceil$, so that $\gamma_n=2^{k'_n}/n^{3/2}\in[2^\theta,2^{\theta+1})$. Every subsequence of $(\gamma_n)$ has a further
subsequence converging to some $\gamma\in[2^\theta,2^{\theta+1}]$, along which $d_n(k'_n)\to L(\gamma)\le L(2^\theta)<\varepsilon$ by Theorem~\ref{thm:bd}; hence $\limsup d_n(k'_n)\le\varepsilon$. If $k_n-\frac32\log_2n\to\infty$,
then $k_n\ge k'_n$ for large $n$, and $d_n(k_n)\le d_n(k'_n)$. So $\limsup d_n(k_n)\le\varepsilon$ for every $\varepsilon$. The other half is the same with $L(2^\theta)>1-\varepsilon$, $\theta$ very negative, and $k'_n=\lfloor\frac32\log_2n+\theta\rfloor$.
\end{proof}

\emph{Why $\frac32$?} The proof shows where the exponent comes from, and it is worth saying in one line. The likelihood ratio is $e^{-nh/a}$ to first order, and $h$ fluctuates
on the scale $\sqrt n$ under the uniform law; so the shuffle is detectable as long as $n\sqrt n/a$ is not small. The $n$ comes from the product of $n$ factors, the $\sqrt n$ from
the central limit theorem, and together they give $a\approx n^{3/2}$, that is, $\frac32\log_2n$ riffles. Compare with random transpositions, where the time
$\frac12n\ln n$ comes from the coupon collector: different shuffles, different statistics, different reasons.

\begin{beauty}
A formula from 1992 for $52!$ probabilities, a carry-counting trick from elementary arithmetic, and the central limit theorem combine into a
single curve, $1-2\Phi(-2^{-\theta}/(4\sqrt3))$, which describes every deck of every size after any number of riffles. And the curve has a meaning: it is
the distance between two normal laws, one shifted. The shuffled deck is a uniform deck with a faint memory of its order, and that memory has
a name, the number of rising sequences.
\end{beauty}

\section{Seven shuffles, and why the answer depends on the question}

On 9 January 1990 the \emph{New York Times} ran a story by Gina Kolata under the headline ``In shuffling cards, 7 is winning number'', reporting
the work of Bayer and Diaconis before its publication. Seven is where the distance for $52$ cards first falls below $\frac12$ ($0.334$); after six it is
still $0.614$. Theorem~\ref{thm:bd} gives $L(2^{7-8.55})\approx0.33$ for seven shuffles: the asymptotic curve is accurate at $n=52$.

But ``seven'' is an answer to a precise question: total variation distance, Gilbert--Shannon--Reeds riffles, from the ordered deck. Change the
question and the answer moves. In 2000 Lloyd N.~Trefethen and Lloyd M.~Trefethen measured randomness by entropy instead, and found that about $\log_2n$
shuffles, six for $52$ cards, already remove all but a small proportion of the information, with about $0.06$ bits left after $\frac32\log_2n$, whatever $n$; a gambler who can see only some features of the deck may need fewer, a magician exploiting rising sequences
needs more to be fooled. The lesson for a reader: \emph{a theorem answers the question it states, and the headline answers a simpler one}. Always
find the definitions before quoting the number.

\section{What the reading taught}

A classic is read differently from a new paper. The first pass is easy, because the claims are famous. The second pass is a map, and the
map is the main product: it tells you what you own (here, every idea of the 1981 paper), what you can sharpen (the constant), and what nobody had
done for decades (the limit profile, Remark~\ref{rem:teyssier}). The third pass is then not a check but an extension: take the next question, and
answer it with the tools the map shows you have. For riffles we could; for transpositions, the profile needs more than we have, and we say so.

\section{Exercises}

\begin{exercise}[$\star$]\label{ex:lemmalimit}
Show that the sum $S$ of Volume~11, Proposition~14.5, satisfies $S\ge(n-1)^2(1-\frac2n)^{2t}$, and deduce that at $t=\frac12n\ln n+cn$ the upper bound lemma cannot give a bound on $d(t)$
better than $\frac12e^{-2c}(1+o(1))$.
\end{exercise}

\begin{exercise}[$\star$]\label{ex:tannysmall}
Check Lemma~\ref{lem:tanny} for $n=3$: compute the law of $\lfloor U_1+U_2+U_3\rfloor$ (volumes of slices of the cube) and the number of permutations of $3$ elements with $0$, $1$, $2$ descents.
\end{exercise}

\begin{exercise}[$\star\star$]\label{ex:shoe}
A casino deals blackjack from a shoe of six decks, $n=312$ cards. Using Theorem~\ref{thm:bd} as an approximation, how many riffle shuffles bring the distance below
$\frac12$? Below $\frac1{10}$? Compare with a single deck.
\end{exercise}

\begin{exercise}[$\star\star$]\label{ex:profile}
Assuming Remark~\ref{rem:teyssier}, compute the limit profile of random transpositions at $c=0$, that is $\|\mathrm{Poisson}(2)-\mathrm{Poisson}(1)\|$, and show that
$\|\mathrm{Poisson}(1+x)-\mathrm{Poisson}(1)\|\sim\frac x{e}$ as $x\to0$. Deduce that the profile decays like $e^{-2c-1}$ as $c\to\infty$. Compare with the riffle profile of Theorem~\ref{thm:bd}: how fast does
it decay as $\theta\to\infty$?
\end{exercise}

\begin{exercise}[$\star\star$]\label{ex:untouched}
For random transpositions, show that the expected number of cards never touched in $t$ swaps is $n(1-\frac1n)^{2t}$, and that at $t=\frac12n\ln n+cn$ it tends to $e^{-2c}$. Explain how this
matches the excess $e^{-2c}$ in Remark~\ref{rem:teyssier}.
\end{exercise}

\begin{exercise}[$\star\star\star$]\label{ex:topinsert}
\emph{Top to random, the profile.} In the top-to-random shuffle (Volume~4, Exercise~8.6) the deck is uniform once the original bottom card has reached the top and been
inserted. Using that the time $\tau$ for this has the law of the coupon collector's time, show that $P(\tau>n\ln n+cn)\to1-e^{-e^{-c}}$. Is this the limit profile of the distance?
(It is an upper bound; Diaconis, Fill and Pitman computed the law of this shuffle exactly in 1992.)
\end{exercise}

\begin{exercise}[$\star\star$]\label{ex:bdmap}
Make a reading map, in the style of the map of this chapter, for the proof of Theorem~\ref{thm:bd}: list each ingredient and the place in the series where it is proved.
Which ingredient would you need to change to treat an $a$-shuffle of a deck that starts in a random order with a known number of rising sequences?
\end{exercise}

\section{Hints}

\begin{hint}{ex:lemmalimit}
All terms of $S$ are non-negative. For the asymptotics, $\ln(1-\frac2n)=-\frac2n-\frac2{n^2}+O(n^{-3})$.
\end{hint}

\begin{hint}{ex:tannysmall}
The density of $S_3$ is piecewise quadratic; by symmetry it is enough to compute $P(S_3<1)=\frac16$.
\end{hint}

\begin{hint}{ex:shoe}
Solve $1-2\Phi(-x)\le\frac12$, that is $x\le0.674$, with $x=\frac{n^{3/2}}{4\sqrt3\cdot2^k}$.
\end{hint}

\begin{hint}{ex:profile}
The two Poisson laws cross once. For small $x$, $p_{1+x}(k)-p_1(k)\approx x\,p_1(k)(k-1)$.
\end{hint}

\begin{hint}{ex:untouched}
A given card is untouched by one swap with probability $(1-\frac1n)^2$.
\end{hint}

\begin{hint}{ex:topinsert}
The coupon collector's time with $n$ coupons, centred at $n\ln n$ and divided by $n$, converges to a Gumbel law.
\end{hint}

\begin{hint}{ex:bdmap}
Only Step~2 uses the starting deck.
\end{hint}

\section{Solutions}

\begin{solution}{ex:lemmalimit}
\idea One term of a sum of non-negative terms is a lower bound for the sum.
\writeup Every term $(f^\lambda)^2s(\lambda)^{2t}$ is non-negative, and the term of $\lambda=(n-1,1)$ has $f^\lambda=n-1$ and $s=1-\frac2n$ (Volume~11, after Proposition~14.5). So $S\ge(n-1)^2(1-\frac2n)^{2t}$ and the bound
the lemma gives, $\frac12\sqrt S$, is at least $\frac12(n-1)(1-\frac2n)^t$. With $t=\frac12n\ln n+cn$,
\[
\ln\Bigl[(n-1)\Bigl(1-\frac2n\Bigr)^t\Bigr]=\ln(n-1)-\Bigl(\frac2n+\frac2{n^2}+O(n^{-3})\Bigr)\Bigl(\frac12n\ln n+cn\Bigr)=-2c+o(1),
\]
so the bound is at least $\frac12e^{-2c}(1+o(1))$: no amount of care with the other terms can beat $e^{-2c}$ by more than a constant factor.
\end{solution}

\begin{solution}{ex:tannysmall}
\idea Compute both sides; the lemma predicts $\frac16,\frac46,\frac16$.
\writeup $P(S_3<1)$ is the volume of the simplex $\{u_i\ge0,\ \sum u_i<1\}$, which is $\frac16$; by the symmetry $u\mapsto1-u$, $P(S_3\ge2)=P(S_3\le1)=\frac16$; so $P(\lfloor S_3\rfloor=1)=\frac46$. The permutations
of $123$: $123$ has no descent; $321$ has two; the other four ($132$, $213$, $231$, $312$) have one. So the counts are $1,4,1$ out of $6$, as predicted. These are the Eulerian numbers.
\end{solution}

\begin{solution}{ex:shoe}
\idea The limit is a function of $x=n^{3/2}/(4\sqrt3\,2^k)$ alone; solve for $k$.
\writeup We need $1-2\Phi(-x)\le\frac12$, that is $\Phi(-x)\ge\frac14$, $x\le0.674$, so $2^k\ge\frac{n^{3/2}}{4\sqrt3\times0.674}=\frac{n^{3/2}}{4.67}$. For $n=312$, $n^{3/2}\approx5511$, so $2^k\ge1180$ and $k=11$. For
distance $\frac1{10}$, $1-2\Phi(-x)\le0.1$ means $x\le0.1257$, $2^k\ge\frac{5511}{0.871}\approx6330$, $k=13$. For one deck, $n^{3/2}\approx375$: $2^k\ge80$, $k=7$ for $\frac12$, and $2^k\ge431$, $k=9$ for $\frac1{10}$. Six decks need four more riffles
than one, since $\frac32\log_26\approx3.9$. (A direct computation verifies these against the exact formula.)
\end{solution}

\begin{solution}{ex:profile}
\idea Two Poisson laws with means $1<\mu$ cross at one point, so the distance is a finite sum.
\writeup With $p_\mu(k)=e^{-\mu}\mu^k/k!$, the ratio $p_\mu(k)/p_1(k)=e^{1-\mu}\mu^k$ increases with $k$, so $p_1>p_\mu$ exactly for $k<k^*=\frac{\mu-1}{\ln\mu}$, and the distance is $\sum_{k<k^*}(p_1(k)-p_\mu(k))$. For
$\mu=2$, $k^*=1.44$, so the distance is $(e^{-1}-e^{-2})+(e^{-1}-2e^{-2})=2e^{-1}-3e^{-2}\approx0.330$. For $\mu=1+x$ with $x\to0$, $k^*=\frac x{\ln(1+x)}\to1$ from above, so the sum is over $k=0,1$. The $k=0$ term is $e^{-1}(1-e^{-x})\sim\frac xe$, and the $k=1$ term is
$e^{-1}(1-(1+x)e^{-x})=O(x^2)$; so the distance is $\sim\frac xe$. With
$x=e^{-2c}$ the profile decays like $e^{-2c-1}$, exponentially in $c$. The riffle profile $1-2\Phi(-\frac{2^{-\theta}}{4\sqrt3})\sim\frac{2\cdot2^{-\theta}}{4\sqrt3\sqrt{2\pi}}$ also decays exponentially, like $2^{-\theta}$; in both cases the tail is linear in the
small parameter, because the distance between two close laws in a one-parameter family is linear in the parameter.
\end{solution}

\begin{solution}{ex:untouched}
\idea Linearity of expectation over the $n$ cards.
\writeup A swap picks two positions independently and uniformly; a given card is untouched if neither choice is its position, with probability $(1-\frac1n)^2$, independently over the
swaps. So the expected number of untouched cards is $n(1-\frac1n)^{2t}$, and $\ln[n(1-\frac1n)^{2t}]=\ln n-(n\ln n+2cn)(\frac1n+O(n^{-2}))\to-2c$. Untouched cards are fixed points, and their number is
approximately Poisson$(e^{-2c})$; the touched cards are roughly uniformly mixed among themselves and contribute about a Poisson$(1)$ number of fixed points, independent of the
rest. The sum of independent Poisson laws is Poisson$(1+e^{-2c})$: that is the heuristic behind Teyssier's theorem, and the precise content of his proof is that nothing else survives.
\end{solution}

\begin{solution}{ex:topinsert}
\idea Use the explicit law of the coupon collector's time as a sum of independent geometric variables, or the inclusion--exclusion formula.
\writeup By Volume~4, Exercise~8.6, $\tau$ has the law of the time to collect $n$ coupons. The probability that a given coupon is missed in $m$ draws is $(1-\frac1n)^m$, and by inclusion--exclusion
\[
P(\tau>m)=\sum_{j\ge1}(-1)^{j+1}\binom nj\Bigl(1-\frac jn\Bigr)^m.
\]
With $m=n\ln n+cn$, $\binom nj(1-\frac jn)^m\to\frac{e^{-jc}}{j!}$ for each $j$, and each term is at most $\frac{n^j}{j!}e^{-jm/n}=\frac{e^{-jc}}{j!}$ in absolute value, which is summable; so by dominated convergence for series,
$P(\tau>m)\to\sum_{j\ge1}(-1)^{j+1}\frac{e^{-jc}}{j!}=1-e^{-e^{-c}}$. Since $\tau$ is a strong stationary time, $d(m)\le P(\tau>m)$ (Volume~4, Proposition~8.18), so $\limsup d(n\ln n+cn)\le1-e^{-e^{-c}}$. Whether this
is the limit profile is a different question, which the bound cannot answer: for that one needs the exact law of the shuffled deck, computed by Diaconis, Fill and Pitman in 1992.
\end{solution}

\begin{solution}{ex:bdmap}
\idea A reading map is a table: ingredient, where it is proved.
\writeup The $a$-shuffle model, Volume~4, Definition~8.22; the formula $Q_a$, Volume~4, Theorem~8.23; composition of shuffles, Volume~4, Proposition~8.24;
total variation as a sum of positive parts, Volume~4, Chapter~8; Tanny's lemma, this chapter; Lindeberg's theorem, Volume~4, Theorem~6.10; the expansion of $\ln(1+x)$, Volume~1;
monotonicity of the distance, Volume~4, Lemma~8.10. If the deck starts in a random order, the law after an $a$-shuffle is no longer a function of the rising sequences of
the final arrangement alone: one composes with the starting law, and Step~1 changes. Step~2 (the law of $r$ under the uniform law) does not.
\end{solution}

\section*{Sources and further reading}
\addcontentsline{toc}{section}{Sources and further reading}

\begin{itemize}
\item \emph{Diaconis.} E.~Klarreich, ``For Persi Diaconis' next magic trick\dots'', \emph{Quanta Magazine}, 14 April 2015 (running away at fourteen with Dai Vernon; Gardner's letter
and admission to Harvard). The MacTutor biography ``Persi Warren Diaconis'' (St Andrews) for City College and the Harvard PhD of 1974. P.~Diaconis and R.~Graham, \emph{Magical
Mathematics}, Princeton University Press, 2011, for the magic.
\item \emph{The papers.} P.~Diaconis and M.~Shahshahani, ``Generating a random permutation with random transpositions'', \emph{Z.\ Wahrsch.\ Verw.\ Gebiete} 57 (1981), 159--179.
D.~Bayer and P.~Diaconis, ``Trailing the dovetail shuffle to its lair'', \emph{Ann.\ Appl.\ Probab.}\ 2 (1992), 294--313. P.~Diaconis, \emph{Group Representations in Probability and
Statistics}, IMS Lecture Notes 11, 1988, Chapter~3, for both shuffles in one place.
\item \emph{The limit profiles.} L.~Teyssier, ``Limit profile for random transpositions'', \emph{Ann.\ Probab.}\ 48 (2020), 2323--2343, arXiv:1905.08514. P.~Diaconis, J.~A.~Fill and J.~Pitman,
``Analysis of top to random shuffles'', \emph{Combin.\ Probab.\ Comput.}\ 1 (1992), 135--155. E.~Nestoridi and S.~Olesker-Taylor, ``Limit profiles for reversible Markov chains'',
\emph{Probab.\ Theory Related Fields} 182 (2022), for later developments.
\item \emph{Descents and uniforms.} S.~Tanny, ``A probabilistic interpretation of Eulerian numbers'', \emph{Duke Math.\ J.}\ 40 (1973), 717--722. R.~P.~Stanley, ``Eulerian partitions of a unit
hypercube'', in \emph{Higher Combinatorics}, Reidel, 1977, for a bijective version.
\item \emph{Seven shuffles in the news, and other measures.} G.~Kolata, ``In shuffling cards, 7 is winning number'', \emph{The New York Times}, 9 January 1990. L.~N.~Trefethen and
L.~M.~Trefethen, ``How many shuffles to randomize a deck of cards?'', \emph{Proc.\ R.\ Soc.\ Lond.\ A} 456 (2000), 2561--2568.
\end{itemize}
